%----------------------------------------------------------------------
\section{Structural Analogies}\label{app:analogies}
%----------------------------------------------------------------------

The following interpretive frameworks provide useful
intuition for the discrimination framework but are not
load-bearing for any theorem in the main text. They are
collected here for readers familiar with these areas.

\subsection{Surreal number framing}

Surreal numbers are constructed as $\{L \mid R\}$, where
$L$ is a set of Surreal numbers less than the number being
defined and $R$ is a set greater. This shares the
$\tau/\sigma$ partition structure: $L$ corresponds to
$\sigma$ (what falls below) and $R$ to $\tau$ (what falls
above), with the number itself defined by the cut between
them.

The ``gap'' cell $(0,0)$ corresponds to Surreal numbers
not yet constructed at a given generation --- the unborn
numbers that are neither in $L$ nor $R$. The ``overlap''
cell $(1,1)$ corresponds to the case where an element
falls on both sides of a cut, indicating that the cut
is too coarse. Each $\kappa$-evolution (each new
discriminator) corresponds to a new generation of Surreal
construction, introducing numbers that fill gaps between
existing ones.

The analogy is suggestive but not exact: Surreal numbers
are totally ordered, while $\kappa$-equivalence classes
are partially ordered. The Surreal construction proceeds
by transfinite induction over a fixed well-ordering,
while $\kappa$-evolution is driven by residues and has
no predetermined sequence.

\subsection{Grothendieck construction}

The Grothendieck construction takes an indexed family of
categories (a functor $F: C^{\mathrm{op}} \to \mathbf{Cat}$)
and assembles them into a single fibered category $\int F$
with a canonical projection to $C$. The discrimination
regime $\kappa$ can be viewed through this lens: each
discriminator $d_i$ defines a fiber (the $\tau/\sigma$
partition of the population under $d_i$), and $\kappa$
as a whole is the total space assembling all fibers.

This perspective is subsumed by the sheaf-theoretic
interpretation (Section \ref{sec:sheaf}), which
identifies $\mathcal{C}(S,\kappa)$ as a sheaf topos on
the site of regimes. The Grothendieck construction
provides the fibration viewpoint; the sheaf topos
provides the covering viewpoint. Both describe the same
structure, but the sheaf interpretation connects to the
toroid geometry and the Fano closure, making it the more
natural home for the paper's purposes.

\subsection{CW-complex interpretation}

The discrimination complex under a regime $\kappa$ of
dimension $n$ can be interpreted as a CW-complex
$X_\kappa$ constructed incrementally:
\begin{itemize}[nosep]
  \item The 0-skeleton $X^0$ consists of the initial
    population (objects prior to any discrimination).
  \item Each $\kappa$-evolution attaching discriminator
    $d_i$ glues a 1-cell (a path separating the elements
    that $d_i$ distinguishes).
  \item Higher cells are generated by wedge products of
    existing discriminators: each non-trivial element of
    $\bigwedge^k V$ for $k \geq 2$ corresponds to a
    $k$-cell witnessing the coherence of $k$
    discriminations.
\end{itemize}

Under this interpretation, a residue corresponds to a
hole in the CW-complex: an element in the kernel of the
XOR signature $\chi_d$ is an element for which no
existing cell provides a discrimination. Consuming the
residue via $\kappa$-evolution is the attachment of a new
cell that fills this hole. Gap residues
($\tau = \sigma = 0$) correspond to lateral cell
attachment (extending the complex to new territory);
overlap residues ($\tau = \sigma = 1$) correspond to
upward cell attachment (introducing a coarser
discrimination).

This language is intuitive but redundant with the chain
complex and homology treatment in Section \ref{sec:boundary},
which captures the same content algebraically without
the topological overhead.

% End of Appendix A
